\documentclass[a4paper,12pt]{article}
\usepackage[utf8]{inputenc}
\usepackage[italian]{babel}
\usepackage{geometry}
\usepackage{setspace}
\usepackage{xcolor}
\usepackage{enumitem}
\usepackage{fancyhdr}
\usepackage{titlesec}

% Font Alegreya
\usepackage[sfdefault]{AlegreyaSans}
\usepackage{Alegreya}
\renewcommand{\rmdefault}{Alegreya-LF}

\geometry{
    margin=2.5cm,
    top=3cm,
    bottom=3cm
}

% Impostazioni spaziatura
\onehalfspacing

% Stile delle sezioni
\titleformat{\section}
{\large\bfseries}{\thesection.}{1em}{}

\titleformat{\subsection}
{\normalsize\bfseries}{\thesubsection}{1em}{}

% Colori
\definecolor{vinaccia}{RGB}{102,51,102}
\definecolor{lightgray}{RGB}{128,128,128}

% Header personalizzato
\pagestyle{fancy}
\fancyhf{}
\fancyhead[R]{\textcolor{black}{\small Volontà Interpretante}}
\fancyhead[L]{\textcolor{black}{\small Alice Cortegiani}}
\fancyfoot[C]{\textcolor{lightgray}{\thepage}}
\renewcommand{\headrulewidth}{0.3pt}
\renewcommand{\footrulewidth}{0pt}

\begin{document}

% FRONTESPIZIO
\thispagestyle{empty}
\begin{center}
    \vspace*{2cm}

    {\Large\textbf{CONSERVATORIO DI MUSICA}}\\[0.3cm]
    {\Large\textbf{``BENEDETTO MARCELLO'' -- VENEZIA}}\\[1.5cm]

    {\large Dottorato di Ricerca AFAM in}\\[0.5cm]
    {\LARGE\textbf{\textcolor{vinaccia}{MUSICA, PERFORMANCE E INNOVAZIONE TECNOLOGICA}}}\\[0.3cm]
    {\large XLI Ciclo -- A.A. 2025/2026}\\[2cm]

    \hrule height 1pt
    \vspace{0.5cm}

    {\Large\textbf{PROGETTO DI RICERCA}}\\[1.5cm]

    \begin{center}
    \textbf{Candidato/a:} Alice Cortegiani\\[1cm]

    \textbf{Titolo provvisorio del progetto:}\\[0.3cm]
    {\large Volontà Interpretante \\ \emph{L'Interprete: messaggero degli dèi}}\\[1.2cm]

    \textbf{Parole chiave} (massimo 5):\\[0.3cm]
    Fenomenologia • Ermeneutica • Prassi Interpretativa • [PAROLA 4] • [PAROLA 5]\\[1cm]
    \end{center}

    \vfill

    {\large Anno Accademico 2025/2026}
\end{center}

\newpage

% INIZIO PROGETTO
\setcounter{page}{1}

\begin{center}
{\textcolor{vinaccia}{\textbf{\large LIMITE: 10.000 caratteri spazi compresi}}}\\
{\small (esclusi frontespizio, bibliografia, ed eventuale apparato illustrativo)}
\end{center}

\vspace{1cm}

% ========== SEZIONE 1 ==========
\section{Descrizione del Progetto}
{\small\textcolor{lightgray}{\textit{Caratteri suggeriti: 3.500--4.000 (35--40\% del totale)}}}\\

... E' la volonta che il mondo abbia un senso.\\
Hèrmes è il messaggero degli dèi, guida le anime dell'Ade e ne è protettore delle vie e dei viandanti. Unisce due dimensioni diverse.\\
Hèrmes porta ai mortali le parole degli dèi, che indicano la loro volontà e la loro spiegazione degli eventi. Quando i mortali ascoltano quelle parole, il senso che per essi il mondo possiede si allarga, si inserisce in un senso più ampio. Annunziando le parole degli dèi, Hèrmes è l'interprete degli eventi del mondo. Si pone tra (inter-) il senso che gli eventi posseggono all'interno di uno sguardo puramente umano e il senso che ad essi conferiscono gli dèi e del quale egli è messaggero.\\
Anche la parola ERMENEUTICA indica dunque, innanzitutto, l'unificazione di due dimensioni.\\ %vedi wikipedia ermeneutica
Ogni interpretazione, anche la più accreditata, può essere riveduta, modificata, sostituita.\\

\emph{Quali sono i mondi che necessitano Hèrmes? Perché Hèrmes?}\\

%creatività e anarchia Agamben?

Il progetto emerge da una preliminare osservazione e interrogazione delle "funzioni" e "ruoli sociali e creativi" dell'essere-musicista: COMPOSITORE-ESECUTORE-PERFORMER-INTERPRETE,  al fine di poter ARMARE le responsabilità di ognuno e collegarli, in relazione di Pòlemos creativo.\\
Il progetto si struttura attraverso le relazioni che intercorrono tra opera,
strumento (il clarinetto) e interprete nel processo creativo di ricerca che
porta alla produzione di nuova musica. Mediante analisi del repertorio che
concede un'esplorazione radicale dello strumento, analisi dei segnali prodotti
dallo strumento con l'uso di tecnologie in grado di descriverne il
comportamento spaziale (che coincide con quello timbrico), una scrittura
dedicata all'esplorazione (dello strumento e della scrittura stessa), il
clarinetto si definisce a luogo di pensiero, di esperienza e consapevolezza,
verso un pensiero creativo che possa definirsi in prassi.
%in continua tendenza verso presupposti utopici, in ascolto.


Massimo Cacciari nel descrivere il tema dell'ascolto in Luigi Nono
\cite{Cacciari1995} disegna lo sfondo di una problematica filosofica che assale
la cultura occidentale nelle relazioni tra scrittura, voce e ascolto, %Problemi di
%ordine generale che si riferiscono al significato culturale generale della
%nascita della scrittura alfabetica che in un lento processo vedono il dominio
%della \emph{visione} sull'\emph{ascolto}
in una

\begin{quote}
  progressiva desomatizzazione della voce. Una perdita dell'udire. L'udire
  non è più una funzione fondamentale del comprendersi. \cite{Cacciari1995}
\end{quote}

Egli sottolinea che in molte pratiche artistiche questo può non essere un
problema, ma non nella musica: la musica non può essere senza ascolto.
La perdita della memoria, dell'ascolto, della memoria dell'ascolto è fatale
per la musica. %Questo accade perché ad ogni ascolto muta il testo stesso.
%La musica non può esistere senza un ascolto vivente, attivo.

\emph{Che cos'è un interprete?}

In questo quadro di sensibilità musicale, l'interprete per primo, cercando di
comprendere, di capire che vorrebbe riuscire a comprendere di più ciò che c'è
\emph{«prima del primo suono e dopo l'ultimo suono»}, si fa testimone della
ricerca musicale nel luogo sonoro in cui opera in grado di
\emph{«articolare, accentuare, dare al canto»}, il testo musicale: a chiarire
anche la distanza dall'esecutore \emph{«l'interprete non legge»} il testo
musicale, \emph{«lo riattiva, lo da al canto»} \cite{Cacciari1995}.\\

\emph{Qual è la memoria dell'interprete? Qual è la letteratura da ri-attivare? Per Quale Musica?}

In Materia e Memoria di Bergson la memoria

\begin{quote}
  significa immaginazione o rappresentazione cosciente del passato in quanto passato
  "in vista" del futuro. E' il futuro che chiama il passato alla "coesistenza"
  con un presente dal quale, per altro, differisce per natura. [...] Noi, afferma Bergson,
  diamo a questa parola "un senso speciale" irriducibile a quello d'uso. Essa vuol dire
  "sintesi del passato e del presente in vista del futuro". [...] Il passato non si risolve
  infatti nel presente come avviene alla confluenza di due fiumi, ma, restando passato
  e proprio in quanto pasato, si media col presente da cui diferisce per natura, per rispondere
  ad un appello proveniente dal futuro. \cite{Ronchi1990}
\end{quote}

La memoria è ermeneutica. è uno \emph{sforzo} di interpretazione. \\
L'interprete contemporaneo raccoglie l'appello proveniente dal futuro, e si fa archeologo [...]
Letteratura come Utopia - Guaccero Luz \\



\subsection{Modalità operative, metodologie e tecnologie d'indagine}
% Inserire descrizione delle metodologie specifiche, tecnologie utilizzate, approcci teorici
% Focus su: artistic research, neuroscienze musicali, tecnologie digitali

\subsection{Tipologia delle fonti}
% Descrivere fonti primarie e secondarie: archivi, manoscritti, registrazioni,
% interviste, letteratura specialistica, database, heritage veneziano

\subsection{Necessità logistiche}
% Indicare eventuali spostamenti per accesso a fonti, archivi, istituzioni, strumenti specifici

\subsection{Difficoltà previste e soluzioni}
% Identificare ostacoli: logistici, tecnici, economici, di accesso alle fonti
% Proporre strategie di risoluzione concrete
% Lorem ipsum dolor sit amet, consectetuer adipiscing elit. Aenean commodo ligula eget dolor. Aenean massa. Cum sociis natoque penatibus et magnis dis parturient montes, nascetur ridiculus mus. Donec quam felis, ultricies nec, pellentesque eu, pretium quis, sem. Nulla consequat massa quis enim. Donec pede justo, fringilla vel, aliquet nec, vulputate eget, arcu. In enim justo, rhoncus ut, imperdiet a, venenatis vitae, justo. Nullam dictum felis eu pede mollis pretium. Integer tincidunt. Cras dapibus. Vivamus elementum semper nisi. Aenean vulputate eleifend tellus. Aenean leo ligula, porttitor eu, consequat vitae, eleifend ac, enim. Aliquam lorem ante, dapibus in, viverra quis, feugiat a, tellus. Phasellus viverra nulla ut metus varius laoreet. Quisque rutrum. Aenean imperdiet. Etiam ultricies nisi vel augue. Curabitur ullamcorper ultricies nisi. Nam eget dui. Etiam rhoncus. Maecenas tempus, tellus eget condimentum rhoncus, sem quam semper libero, sit amet adipiscing sem neque sed ipsum. Nam quam nunc, blandit vel, luctus pulvinar, hendrerit id, lorem. Maecenas nec odio et ante tincidunt tempus. Donec vitae sapien ut libero venenatis faucibus. Nullam quis ante. Etiam sit amet orci eget eros faucibus tincidunt. Duis leo. Sed fringilla mauris sit amet nibh. Donec sodales sagittis magna. Sed consequat, leo eget bibendum sodales, augue velit cursus nunc, quis gravida magna mi a libero. Fusce vulputate eleifend sapien. Vestibulum purus quam, scelerisque ut, mollis sed, nonummy id, metus. Nullam accumsan lorem in dui. Cras ultricies mi eu turpis hendrerit fringilla. Vestibulum ante ipsum primis in faucibus orci luctus et ultrices posuere cubilia Curae; In ac dui quis mi consectetuer lacinia. Nam pretium turpis et arcu. Duis arcu tortor, suscipit eget, imperdiet nec, imperdiet iaculis, ipsum. Sed aliquam ultrices mauris. Integer ante arcu, accumsan a, consectetuer eget, posuere ut, mauris. Praesent adipiscing. Phasellus ullamcorper ipsum rutrum nunc. Nunc nonummy metus. Vestibulum volutpat pretium libero. Cras id dui. Aenean ut eros et nisl sagittis vestibulum. Nullam nulla eros, ultricies sit amet, nonummy id, imperdiet feugiat, pede. Sed lectus. Donec mollis hendrerit risus. Phasellus nec sem in justo pellentesque facilisis. Etiam imperdiet imperdiet orci. Nunc nec neque. Phasellus leo dolor, tempus non, auctor et, hendrerit quis, nisi. Curabitur ligula sapien, tincidunt non, euismod vitae, posuere imperdiet, leo. Maecenas malesuada. Praesent congue erat at massa. Sed cursus turpis vitae tortor. Donec posuere vulputate arcu. Phasellus accumsan cursus velit. Vestibulum ante ipsum primis in faucibus orci luctus et ultrices posuere cubilia Curae; Sed aliquam, nisi quis porttitor congue, elit erat euismod orci, ac placerat dolor lectus quis orci. Phasellus consectetuer vestibulum elit. Aenean tellus metus, bibendum sed, posuere ac, mattis non, nunc. Vestibulum fringilla pede sit amet augue. In turpis. Pellentesque posuere. Praesent turpis. Aenean posuere, tortor sed cursus feugiat, nunc augue blandit nunc, eu sollicitudin urna dolor sagittis lacus. Donec elit libero, sodales nec, volutpat a, suscipit non, turpis. Nullam sagittis. Suspendisse pulvinar, augue ac venenatis condimentum, sem libero volutpat nibh, nec pellentesque velit pede quis nunc. Vestibulum ante ipsum primis in faucibus orci luctus et ultrices posuere cubilia Curae; Fusce id purus. Ut varius tincidunt libero. Phasellus dolor. Maecenas vestibulum mollis diam. Pellentesque ut neque. P

% ========== SEZIONE 2 ==========
\section{Stato dell'Arte}
{\small\textcolor{lightgray}{\textit{Caratteri suggeriti: 2.000--2.500 (20--25\% del totale)}}}\\

Il progetto nasce da un'attività quotidiana di ricerca presso laboratori
e centri specializzati, in una pratica musicale attenta alle necessità di un
pensare creativo e analitico, dove il ruolo dell'interprete possa rinascere
dalle ceneri dell'intrattenimento cameristico, operistico e sinfonico.\\
HERITAGE VENEZIANO

% Inquadrare il progetto nel panorama degli studi esistenti
% Indicare se si parte da esperienze pregresse o percorso nuovo
% Citare contributi rilevanti in: performance studies, neuroscienze della musica,
% tecnologie musicali, didattica speciale, heritage culturale
% Lorem ipsum dolor sit amet, consectetuer adipiscing elit. Aenean commodo ligula eget dolor. Aenean massa. Cum sociis natoque penatibus et magnis dis parturient montes, nascetur ridiculus mus. Donec quam felis, ultricies nec, pellentesque eu, pretium quis, sem. Nulla consequat massa quis enim.
Donec pede justo, fringilla vel, aliquet nec, vulputate eget, arcu. In enim justo, rhoncus ut, imperdiet a, venenatis vitae, justo. Nullam dictum felis eu pede mollis pretium. Integer tincidunt. Cras dapibus. Vivamus elementum semper nisi. Aenean vulputate eleifend tellus. Aenean leo ligula, porttitor eu, consequat vitae, eleifend ac, enim. Aliquam lorem ante, dapibus in, viverra quis, feugiat a, tellus. Phasellus viverra nulla ut metus varius laoreet. Quisque rutrum. Aenean imperdiet. Etiam ultricies nisi vel augue. Curabitur ullamcorper ultricies nisi. Nam eget dui. Etiam rhoncus. Maecenas tempus, tellus eget condimentum rhoncus, sem quam semper libero, sit amet adipiscing sem neque sed ipsum. Nam quam nunc, blandit vel, luctus pulvinar, hendrerit id, lorem. Maecenas nec odio et ante tincidunt tempus. Donec vitae sapien ut libero venenatis faucibus. Nullam quis ante. Etiam sit amet orci eget eros faucibus tincidunt. Duis leo. Sed fringilla mauris sit amet nibh. Donec sodales sagittis magna. Sed consequat, leo eget bibendum sodales, augue velit cursus nunc, quis gravida magna mi a libero. Fusce vulputate eleifend sapien. Vestibulum purus quam, scelerisque ut, mollis sed, nonummy id, metus. Nullam accumsan lorem in dui. Cras ultricies mi eu turpis hendrerit fringilla. Vestibulum ante ipsum primis in faucibus orci luctus et ultrices posuere cubilia Curae; In ac dui quis mi consectetuer lacinia. Nam pretium turpis et arcu. Duis arcu tortor, suscipit eget, imperdiet nec, imperdiet iaculis, ipsum. Sed aliquam ultrices mauris. Integer ante arcu, accumsan a, consectetuer eget, posuere ut, mauris. Praesent adipiscing. Phasellus ullamcorper ipsum rutrum nunc. Nunc nonummy metus. Vestibulum volutpat pretium libero. Cras id dui. Aenean ut eros et nisl sagittis vestibulum. Nullam nulla eros, ultricies sit amet, nonummy id, imperdiet feugiat, pede. Sed lectus. Donec mollis hendrerit risus. Phasellus nec sem in justo pellentesque facilisis. Etiam imperdiet imperdiet orci. Nunc nec neque. Phasellus leo dolor, tempus non, auctor et, hendrerit quis, nisi. Curabitur ligula sapien, tincidunt non, euismod vitae, posuere imperdiet, leo. Maecenas malesuada. Praesent congue erat at massa..

\subsection{Esperienza pregressa del candidato}

\subsection{Posizionamento innovativo}
% ========== SEZIONE 3 ==========
\section{Risultati Attesi}
{\small\textcolor{lightgray}{\textit{Caratteri suggeriti: 2.000--2.500 (20--25\% del totale)}}}\\

\subsection{Innovazioni rispetto allo stato dell'arte attuale}
% Specificare avanzamenti rispetto alle conoscenze attuali
% Evidenziare originalità nel dialogo arte-scienza-tecnologie

\subsection{Competenze e punti di forza del candidato}
% Illustrare idoneità per questa ricerca: background, esperienze,
% competenze performative, tecniche, di ricerca

\subsection{Impatto previso}

Lorem ipsum dolor sit amet, consectetuer adipiscing elit. Aenean commodo ligula eget dolor. Aenean massa. Cum sociis natoque penatibus et magnis dis parturient montes, nascetur ridiculus mus. Donec quam felis, ultricies nec, pellentesque eu, pretium quis, sem. Nulla consequat massa quis enim. Donec pede justo, fringilla vel, aliquet nec, vulputate eget, arcu. In enim justo, rhoncus ut, imperdiet a, venenatis vitae, justo. Nullam dictum felis eu pede mollis pretium. Integer tincidunt. Cras dapibus. Vivamus elementum semper nisi. Aenean vulputate eleifend tellus. Aenean leo ligula, porttitor eu, consequat vitae, eleifend ac, enim. Aliquam lorem ante, dapibus in, viverra quis, feugiat a, tellus. Phasellus viverra nulla ut metus varius laoreet. Quisque rutrum. Aenean imperdiet. Etiam ultricies nisi vel augue. Curabitur ullamcorper ultricies nisi. Nam eget dui. Etiam rhoncus. Maecenas tempus, tellus eget condimentum rhoncus, sem quam semper libero, sit amet adipiscing sem neque sed ipsum. Nam quam nunc, blandit vel, luctus pulvinar, hendrerit id, lorem. Maecenas nec odio et ante tincidunt tempus. Donec vitae sapien ut libero venenatis faucibus. Nullam quis ante. Etiam sit amet orci eget eros faucibus tincidunt. Duis leo. Sed fringilla mauris sit amet nibh. Donec sodales sagittis magna. Sed consequat, leo eget bibendum sodales, augue velit cursus nunc, quis gravida magna mi a libero. Fusce vulputate eleifend sapien. Vestibulum purus quam, scelerisque ut, mollis sed, nonummy id, metus. Nullam accumsan lorem in dui. Cras ultricies mi eu turpis hendrerit fringilla. Vestibulum ante ipsum primis in faucibus orci luctus et ultrices posuere cubilia Curae; In ac dui quis mi consectetuer lacinia. Nam pretium turpis et arcu. Duis arcu tortor, suscipit eget, imperdiet nec, imperdiet iaculis, ipsum. Sed aliquam ultrices mauris. Integer ante arcu, accumsan a, consectetuer eget, posuere ut, mauris. Praesent adipiscing. Phasellus ullamcorper ipsum rutrum nunc. Nunc nonummy metus. Vestibulum volutpat pretium libero. Cras id dui. Aenean ut eros et nisl sagittis vestibulum. Nullam nulla eros, ultricies sit amet, nonummy id, imperdiet feugiat, pede. Sed lectus. Donec mollis hendrerit risus. Phasellus nec sem in justo pellentesque facilisis. Etiam imperdiet imperdiet orci. Nunc nec neque. Phasellus leo dolor, tempus non, auctor et, hendrerit quis, nisi. Curabitur ligula sapien, tincidunt non, euismod vitae, posuere imperdiet, leo. Maecenas malesuada. Praesent congue erat at massa..

% ========== SEZIONE 4 ==========
\section{Distribuzione del Lavoro nei Tre Anni}
{\small\textcolor{lightgray}{\textit{Caratteri suggeriti: 1.500--2.000 (15--20\% del totale)}}}\\

\textbf{Primo Anno:}
% Attività preparatorie, corsi, seminari, ricerca bibliografica, primi approcci alle fonti

\textbf{Secondo Anno:}
% Ricerca principale, raccolta dati, analisi, viaggi di studio, performance sperimentali

\textbf{Terzo Anno:}
% Elaborazione risultati, scrittura tesi, performance finali, disseminazione risultati
Lorem ipsum dolor sit amet, consectetuer adipiscing elit. Aenean commodo ligula eget dolor. Aenean massa. Cum sociis natoque penatibus et magnis dis parturient montes, nascetur ridiculus mus. Donec quam felis, ultricies nec, pellentesque eu, pretium quis, sem. Nulla consequat massa quis enim. Donec pede justo, fringilla vel, aliquet nec, vulputate eget, arcu. In enim justo, rhoncus ut, imperdiet a, venenatis vitae, justo. Nullam dictum felis eu pede mollis pretium. Integer tincidunt. Cras dapibus. Vivamus elementum semper nisi. Aenean vulputate eleifend tellus. Aenean leo ligula, porttitor eu, consequat vitae, eleifend ac, enim. Aliquam lorem ante, dapibus in, viverra quis, feugiat a, tellus. Phasellus viverra nulla ut metus varius laoreet. Quisque rutrum. Aenean imperdiet. Etiam ultricies nisi vel augue. Curabitur ullamcorper ultricies nisi. Nam eget dui. Etiam rhoncus. Maecenas tempus, tellus eget condimentum rhoncus, sem quam semper libero, sit amet adipiscing sem neque sed ipsum. Nam quam nunc, blandit vel, luctus pulvinar, hendrerit id, lorem. Maecenas nec odio et ante tincidunt tempus. Donec vitae sapien ut libero venenatis faucibus. Nullam quis ante. Etiam sit amet orci eget eros faucibus tincidunt. Duis leo. Sed fringilla mauris sit amet nibh. Donec sodales sagittis magna. Sed consequat, leo eget bibendum sodales, augue velit cursus nunc, quis gravida magna mi a libero. Fusce vulputate eleifend sapien. Vestibulum purus quam, scelerisque ut, mollis sed, nonummy id, met


\newpage

% ========== BIBLIOGRAFIA ==========
\section*{Bibliografia Iniziale}
{\small\textcolor{lightgray}{\textit{Non concorre al conteggio dei caratteri}}}

\begin{itemize}[leftmargin=1cm, itemsep=0.5em]
    \item[] % [Inserire qui le principali fonti bibliografiche di riferimento]

    \item[] % Esempio formato:
    % Cognome, N. (Anno). \textit{Titolo}. Casa Editrice.

    \item[]

    \item[]

    \item[]

    \item[]

    \item[]

    \item[]

    \item[]

    \item[]
\end{itemize}

\vspace{2cm}

% ========== ISTRUZIONI FINALI ==========
\begin{center}
\fbox{
\begin{minipage}{0.9\textwidth}
\textbf{\textcolor{vinaccia}{PROMEMORIA FINALE}}

\begin{itemize}[leftmargin=0.5cm]
    \item \textbf{Controllo caratteri:} Verificare che il testo principale non superi i 10.000 caratteri
    \item \textbf{Tematiche chiave:} Il progetto deve coniugare competenze performative con artistic research
    \item \textbf{Focus suggeriti:} Arte-scienza-tecnologie, heritage veneziano, neuroscienze musicali, didattica speciale per l'inclusione
    \item \textbf{Originalità:} Evidenziare gli aspetti originali e l'impatto atteso della ricerca
    \item \textbf{Fattibilità:} Dimostrare che il progetto è realisticamente realizzabile in tre anni
\end{itemize}
\end{minipage}
}
\end{center}

\end{document}
